\chapter{Localisation and Dimension}

\begin{center}\itshape
To study a ring near one of its primes, invert everything outside it. Localisation makes this precise, the prime spectrum turns every ring into a space, and integral extensions
measure how many independent directions that space has: its dimension.
\end{center}

To decide whether an integer is a square, it helps to look at it one prime at a time: modulo $8$, modulo $3$, modulo $5$. To study a curve, one looks at it near each of its points. Both
habits have the same algebraic form. Near the prime $p$, every integer prime to $p$ behaves like a unit, so one makes it a unit: the ring $\Z_{(p)}$ of fractions $a/b$ with $p\nmid b$
sees only the prime $p$. Near the point $a$ of a variety, a polynomial that does not vanish at $a$ can be inverted, and the ring of fractions $f/g$ with $g(a)\ne0$ sees only what happens close to $a$.
Wolfgang Krull, who had worked with Emmy Noether in Göttingen, turned these \emph{local rings} into the main tool of commutative algebra in the 1930s, and the operation that produces them is called
\emph{localisation}.

A second idea of the period is that of dimension. A plane curve is one-dimensional, a surface two-dimensional, and Krull defined the dimension of any ring as the length of the
longest chain of prime ideals: in $\C[x,y]$, $0\subsetneq(y)\subsetneq(x,y)$, a surface containing a curve containing a point. That this gives $\dim K[x_1,\dots,x_n]=n$ is not obvious at
all, and the proof uses Emmy Noether's normalisation lemma of 1926.

In the 1960s Alexander Grothendieck took the last step: every ring is the ring of functions on a space, its \emph{spectrum}, whose points are its prime ideals. David Mumford's notes on
algebraic geometry contain a famous drawing of the spectrum of $\Z[x]$, a two-dimensional ``arithmetic surface'' in which the primes of $\Z$ are vertical lines and the polynomials are
curves. Figure~\ref{fig:specz} is a version of it, and this chapter proves that it is accurate.

All rings are commutative with $1$.

\section{Rings of fractions}

A subset $S\subseteq R$ is \emph{multiplicative} if $1\in S$ and $st\in S$ whenever $s,t\in S$. Examples: the powers $\{1,f,f^2,\dots\}$ of an element; the complement $R\setminus\mathfrak p$ of a prime ideal; the
non-zero elements of a domain.

\begin{definition}\label{def:localisation}
For a multiplicative $S$, the \emph{ring of fractions} $S^{-1}R$ is the set of pairs $(r,s)\in R\times S$ modulo the relation $(r,s)\sim(r',s')$ if $t(rs'-r's)=0$ for some $t\in S$, with the class of $(r,s)$ written $r/s$ and
the operations $\frac rs+\frac{r'}{s'}=\frac{rs'+r's}{ss'}$, $\frac rs\cdot\frac{r'}{s'}=\frac{rr'}{ss'}$. The map $\iota\colon R\to S^{-1}R$, $r\mapsto r/1$, is a ring homomorphism. For $S=\{f^n\}$ one writes $R_f$, and for $S=R\setminus\mathfrak p$, $R_{\mathfrak p}$,
the \emph{localisation of $R$ at $\mathfrak p$}. For a module $M$, $S^{-1}M$ is defined in the same way, with $m/s$ for $m\in M$.
\end{definition}

The factor $t$ is needed when $R$ has zero divisors: without it the relation would not be transitive. Checking that $\sim$ is an equivalence relation and that the operations are well defined is routine: for
transitivity, $t(rs'-r's)=0$ and $t'(r's''-r''s')=0$ give $tt's'(rs''-r''s)=0$, after multiplying the first by $t's''$, the second by $ts$, and adding. If $R$ is a domain and $0\notin S$, the factor
$t$ can be cancelled, and $S^{-1}R$ is the subring of the field of fractions made of the $r/s$. So $\Z_{(p)}=\{a/b:p\nmid b\}$, $\Z[1/6]=\Z_6$ is the ring of fractions with denominators $2^a3^b$, and the
field of fractions of a domain is its localisation at the prime $(0)$. If $0\in S$, then $S^{-1}R=0$.

\begin{proposition}[universal property]\label{prop:locuniv}
$\iota(s)$ is a unit of $S^{-1}R$ for every $s\in S$, and every ring homomorphism $\varphi\colon R\to T$ such that $\varphi(s)$ is a unit for every $s\in S$ factors uniquely as $\varphi=\bar\varphi\circ\iota$. Moreover
$\iota(r)=0$ if and only if $sr=0$ for some $s\in S$.
\end{proposition}

\begin{proof}
$\iota(s)$ has inverse $1/s$. The only possible $\bar\varphi$ is $\bar\varphi(r/s)=\varphi(r)\varphi(s)^{-1}$, and it is well defined: if $t(rs'-r's)=0$, then $\varphi(t)$ is a unit, so $\varphi(r)\varphi(s')=\varphi(r')\varphi(s)$. It is a homomorphism by the
same computation as for fractions. Finally $r/1=0/1$ means $t(r\cdot1-0)=0$ for some $t\in S$.
\end{proof}

So $S^{-1}R$ is the universal way of making $S$ invertible, and Chapter~1's language says it once more: $R\mapsto S^{-1}R$ is left adjoint to the inclusion of the $R$-algebras in which $S$ becomes invertible.

\begin{proposition}[localisation is exact]\label{prop:locexact}
\leavevmode
\begin{enumerate}[label=(\alph*)]
\item There is a natural isomorphism $S^{-1}R\otimes_RM\cong S^{-1}M$, $\frac rs\otimes m\mapsto\frac{rm}s$.
\item If $M'\xrightarrow fM\xrightarrow gM''$ is exact, so is $S^{-1}M'\to S^{-1}M\to S^{-1}M''$, with the maps $m/s\mapsto f(m)/s$ and $m/s\mapsto g(m)/s$.
\item Hence $S^{-1}R$ is a flat $R$-module. In particular $\Q$ is a flat $\Z$-module.
\end{enumerate}
\end{proposition}

\begin{proof}
(b) The composite is $0$ as $g\circ f=0$. If $g(m)/s=0$ in $S^{-1}M''$, then $tg(m)=g(tm)=0$ for some $t\in S$, so $tm=f(m')$ for some $m'$, and $m/s=f(m')/(ts)$ is in the image.
(a) The map is induced by the bilinear map $(r/s,m)\mapsto rm/s$, which is well defined. It is surjective, since $m/s$ is the image of $\frac1s\otimes m$. For injectivity, every element of $S^{-1}R\otimes M$ can be written
$\frac1s\otimes m$: bring a sum $\sum\frac{r_i}{s_i}\otimes m_i$ to the common denominator $s=\prod s_i$, then $\frac{a_i}s\otimes m_i=\frac1s\otimes a_im_i$. If $m/s=0$ in $S^{-1}M$, then $tm=0$ for some $t\in S$, and
$\frac1s\otimes m=\frac t{ts}\otimes m=\frac1{ts}\otimes tm=0$. Naturality is clear on generators.
(c) By (a) and (b), $S^{-1}R\otimes-$ preserves injective maps. For $\Q=S^{-1}\Z$ with $S=\Z\setminus\{0\}$, this keeps the promise of Chapter~2.
\end{proof}

The ideals of a ring of fractions are easy to describe, and its primes are exactly what one hopes.

\begin{theorem}[primes of a ring of fractions]\label{thm:locprimes}
The maps $\mathfrak p\mapsto S^{-1}\mathfrak p=\{p/s:p\in\mathfrak p,s\in S\}$ and $\mathfrak q\mapsto\iota^{-1}(\mathfrak q)$ are inverse, inclusion-preserving bijections between the prime ideals of $R$ that do not meet $S$ and the prime
ideals of $S^{-1}R$.
\end{theorem}

\begin{proof}
For a prime $\mathfrak q$ of $S^{-1}R$, $\iota^{-1}(\mathfrak q)$ is a prime of $R$ (an inverse image of a prime ideal), and it does not meet $S$, since $\iota(s)$ is a unit. Moreover $S^{-1}(\iota^{-1}(\mathfrak q))=\mathfrak q$: if $r/s\in\mathfrak q$
then $r/1=s\cdot r/s\in\mathfrak q$, so $r\in\iota^{-1}(\mathfrak q)$. Conversely let $\mathfrak p$ be prime with $\mathfrak p\cap S=\varnothing$. Then $S^{-1}\mathfrak p$ is an ideal. If $r/1\in S^{-1}\mathfrak p$, then $r/1=p/s$, so $t(rs-p)=0$ for some
$t\in S$, and $tsr=tp\in\mathfrak p$ with $ts\notin\mathfrak p$, hence $r\in\mathfrak p$: so $\iota^{-1}(S^{-1}\mathfrak p)=\mathfrak p$, which also shows $1/1\notin S^{-1}\mathfrak p$. Finally $S^{-1}\mathfrak p$ is prime: if
$\frac rs\frac{r'}{s'}\in S^{-1}\mathfrak p$, then $rr'/1$ is in it, so $rr'\in\mathfrak p$, and $r$ or $r'$ is in $\mathfrak p$.
\end{proof}

\begin{corollary}\label{cor:localring}
For a prime $\mathfrak p$, the ring $R_{\mathfrak p}$ has exactly one maximal ideal, $\mathfrak pR_{\mathfrak p}$, and its primes correspond to the primes of $R$ contained in $\mathfrak p$.
\end{corollary}

\begin{proof}
The primes of $R$ not meeting $R\setminus\mathfrak p$ are those contained in $\mathfrak p$, and among them $\mathfrak p$ is the largest. Every proper ideal $I$ of $R_{\mathfrak p}$ lies in $\mathfrak pR_{\mathfrak p}$: if $r/s\in I$ with $r\notin\mathfrak p$, then $r/s$ is
a unit with inverse $s/r$, and $I$ is not proper.
\end{proof}

\section{Local rings and Nakayama's lemma}

A ring with exactly one maximal ideal $\mathfrak m$ is a \emph{local ring}; its \emph{residue field} is $R/\mathfrak m$. In a local ring the units are exactly the elements outside $\mathfrak m$: a non-unit $r$ generates a proper ideal,
which lies in a maximal ideal, that is in $\mathfrak m$ (for this one uses Zorn's lemma, Volume~6, Theorem~1.19, or, in a Noetherian ring, Proposition~\ref{prop:noetherian}). Examples: fields, $\Z_{(p)}$ with residue field $\mathbb F_p$, the
formal power series $K[[x]]$, and the ring of germs of smooth functions at a point.

\begin{theorem}[Nakayama's lemma]\label{thm:nakayama}
Let $(R,\mathfrak m)$ be a local ring and $M$ a finitely generated $R$-module with $\mathfrak mM=M$. Then $M=0$.
\end{theorem}

\begin{proof}
Suppose $M\ne0$ and let $m_1,\dots,m_n$ be a generating family with $n\ge1$ minimal. Since $m_n\in\mathfrak mM$, $m_n=\sum_ia_im_i$ with $a_i\in\mathfrak m$, so $(1-a_n)m_n=\sum_{i<n}a_im_i$. As $a_n\in\mathfrak m$, $1-a_n$ is a unit, and
$m_n$ is a combination of $m_1,\dots,m_{n-1}$: these generate $M$, contradicting minimality (if $n=1$, the conclusion is $m_1=0$ and $M=0$).
\end{proof}

\begin{corollary}\label{cor:mingens}
Let $(R,\mathfrak m)$ be local with residue field $k$ and $M$ finitely generated. Elements $m_1,\dots,m_n\in M$ generate $M$ if and only if their images span the $k$-vector space $M/\mathfrak mM$. So every minimal generating family of
$M$ has $\dim_kM/\mathfrak mM$ elements.
\end{corollary}

\begin{proof}
If the images span, let $N=\sum Rm_i$. Then $M=N+\mathfrak mM$, so $\mathfrak m(M/N)=M/N$, and $M/N=0$ by Theorem~\ref{thm:nakayama}. The converse is clear. A generating family is minimal if and only if its images form a basis.
\end{proof}

Nakayama's lemma fails without finite generation: $\Q$, as a module over $\Z_{(p)}$, satisfies $p\Q=\Q$. Its typical use is to pass from information modulo $\mathfrak m$, which is linear algebra over a field, to
information about $M$. It also shows that being zero is a local property.

\begin{proposition}\label{prop:localzero}
For a module $M$, the following are equivalent: $M=0$; $M_{\mathfrak p}=0$ for every prime $\mathfrak p$; $M_{\mathfrak m}=0$ for every maximal ideal $\mathfrak m$. Consequently a homomorphism is injective, or surjective, if and only if
it is so after localisation at every maximal ideal.
\end{proposition}

\begin{proof}
Suppose $M_{\mathfrak m}=0$ for all maximal $\mathfrak m$, and let $x\in M$. Its annihilator $\operatorname{Ann}(x)=\{r:rx=0\}$ is an ideal. If it were proper, it would lie in a maximal ideal $\mathfrak m$; but $x/1=0$ in $M_{\mathfrak m}$ means
$sx=0$ for some $s\notin\mathfrak m$, so $s\in\operatorname{Ann}(x)\setminus\mathfrak m$, a contradiction. So $1\cdot x=0$. For a homomorphism $u$, apply this to $\ker u$ and $\operatorname{coker}u$, which commute with localisation by
Proposition~\ref{prop:locexact}(b).
\end{proof}

\section{The spectrum of a ring}

\begin{definition}\label{def:spec}
The \emph{spectrum} $\operatorname{Spec}R$ is the set of prime ideals of $R$. For an ideal $I$, let $V(I)=\{\mathfrak p:\mathfrak p\supseteq I\}$. These sets are the closed sets of a topology, the \emph{Zariski topology} of $\operatorname{Spec}R$. A ring
homomorphism $\varphi\colon R\to T$ induces the map $\operatorname{Spec}T\to\operatorname{Spec}R$, $\mathfrak q\mapsto\varphi^{-1}(\mathfrak q)$, which is continuous.
\end{definition}

The axioms are checked as in Chapter~3, Proposition~\ref{prop:zariski}: $V(I)\cup V(J)=V(IJ)$ because a prime containing $IJ$ contains $I$ or $J$, and $\bigcap V(I_\lambda)=V(\sum I_\lambda)$. Continuity holds because the
inverse image of $V(I)$ is $V(\varphi(I)T)$. So $\operatorname{Spec}$ is a contravariant functor from rings to topological spaces. The closed points are the maximal ideals, since the closure of $\{\mathfrak p\}$ is $V(\mathfrak p)$.
Over an algebraically closed field $K$, the maximal ideals of $K[x_1,\dots,x_n]$ are the points of $K^n$ (Theorem~\ref{thm:weaknull}), so $\operatorname{Spec}K[x_1,\dots,x_n]$ contains $K^n$, and one more point for each
irreducible closed subset: its prime ideal, a \emph{generic point} whose closure is the whole subset.

$\operatorname{Spec}\Z$ consists of $(0)$ and the $(p)$. Its closed sets are the finite sets of non-zero primes and the whole space, and the generic point $(0)$ is in every non-empty open set: a line with a point that is everywhere.
The next example is the one Mumford drew.

\begin{theorem}[the primes of $\Z\lbrack x\rbrack$]\label{thm:speczx}
The prime ideals of $\Z[x]$ are:
\begin{enumerate}[label=(\alph*)]
\item $(0)$;
\item the principal ideals $(f)$, for $f$ irreducible in $\Z[x]$: a prime number $p$, or a primitive polynomial irreducible over $\Q$;
\item the maximal ideals $(p,f)$, for $p$ prime and $f\in\Z[x]$ whose reduction modulo $p$ is irreducible in $\mathbb F_p[x]$.
\end{enumerate}
\end{theorem}

\begin{proof}
These ideals are prime: $\Z[x]$ is a unique factorisation domain (Volume~2, Theorem~6.17), so irreducible elements generate prime ideals, and $\Z[x]/(p,f)\cong\mathbb F_p[x]/(\bar f)$ is a field. Conversely let $P$ be prime. Then
$P\cap\Z$ is a prime ideal of $\Z$, so $(0)$ or $(p)$.

If $P\cap\Z=(0)$, then $P$ does not meet $S=\Z\setminus\{0\}$, and by Theorem~\ref{thm:locprimes} it corresponds to a prime of $S^{-1}\Z[x]=\Q[x]$, a principal ideal domain: $(0)$, or $(g)$ with $g$ irreducible in $\Q[x]$. In the
first case $P=(0)$. In the second, scale $g$ to a primitive $f\in\Z[x]$; then $P=f\Q[x]\cap\Z[x]$, which equals $f\Z[x]$ by Gauss's lemma: if $fh\in\Z[x]$ with $h\in\Q[x]$, write $h=ch_0$ with $h_0$ primitive and $c\in\Q$; the
product $fh_0$ is primitive, so $c\in\Z$.

If $P\cap\Z=(p)$, then $P/(p)$ is a prime of $\Z[x]/(p)=\mathbb F_p[x]$, a principal ideal domain: either $(0)$, giving $P=(p)$, or $(\bar f)$ with $\bar f$ irreducible, giving $P=(p,f)$ for any lift $f$ of $\bar f$.
\end{proof}

\begin{figure}[ht]
\centering
\fig{ch04-speczx}
\caption{After Mumford: the spectrum of $\Z[x]$. The bottom line is $\operatorname{Spec}\Z$. Above each prime $p$ stands the fibre $\operatorname{Spec}\mathbb F_p[x]$, a vertical line made of the generic point $(p)$ and the
closed points $(p,f)$, one for each monic irreducible $\bar f$ modulo $p$, of which a few are drawn in grey. The curve is the closure of the point $(x^2+1)$: it meets the fibre at $p$ in two points
$(p,x\pm a)$ when $p\equiv1\pmod4$, in one point $(p,x^2+1)$ of degree $2$ when $p\equiv3\pmod4$, and has its two branches meeting at the single point $(2,x+1)$ of the fibre at $2$, where $x^2+1=(x+1)^2$.}
\label{fig:specz}
\end{figure}

\section{Integral extensions}

\begin{definition}\label{def:integral}
Let $A\subseteq B$ be rings. An element $b\in B$ is \emph{integral} over $A$ if it is a root of a monic polynomial with coefficients in $A$. The ring $B$ is \emph{integral} over $A$ if all its elements are.
\end{definition}

The algebraic integers of Volume~12 are the complex numbers integral over $\Z$. The criterion that makes integrality usable is the ``determinant trick''.

\begin{proposition}\label{prop:integralcrit}
For $b\in B$, the following are equivalent: (a) $b$ is integral over $A$; (b) $A[b]$ is a finitely generated $A$-module; (c) $b$ lies in a subring $C\subseteq B$ containing $A$ that is a finitely generated $A$-module.
\end{proposition}

\begin{proof}
(a)$\Rightarrow$(b): if $b^n=-(a_{n-1}b^{n-1}+\dots+a_0)$, every power of $b$ is an $A$-combination of $1,b,\dots,b^{n-1}$, by induction. (b)$\Rightarrow$(c): take $C=A[b]$. (c)$\Rightarrow$(a): let $C=Ac_1+\dots+Ac_n$. Then
$bc_i=\sum_ja_{ij}c_j$, that is $(bI-(a_{ij}))c=0$ for the column $c=(c_j)$. Multiplying by the adjugate matrix gives $\det(bI-(a_{ij}))c_j=0$ for all $j$, and since $1\in C$ is a combination of the $c_j$, $\det(bI-(a_{ij}))=0$. Expanded, this is a
monic equation for $b$ over $A$.
\end{proof}

\begin{corollary}\label{cor:integralclosure}
The elements of $B$ integral over $A$ form a subring, the \emph{integral closure} of $A$ in $B$. If $C$ is integral over $B$ and $B$ integral over $A$, then $C$ is integral over $A$.
\end{corollary}

\begin{proof}
If $b,b'$ are integral, $A[b,b']$ is a finitely generated $A$-module (generated by the $b^ib'^j$ with $i,j$ below the degrees of the equations), and it contains $b\pm b'$ and $bb'$. For transitivity, $c\in C$ satisfies $c^n+b_{n-1}c^{n-1}+\dots+b_0=0$
with $b_i\in B$; then $A[b_0,\dots,b_{n-1},c]$ is a finitely generated $A$-module, by the first argument repeated, and Proposition~\ref{prop:integralcrit} applies.
\end{proof}

A domain is \emph{integrally closed}, or \emph{normal}, if it equals its integral closure in its field of fractions. Unique factorisation domains are integrally closed (the argument of the rational root theorem), but
$\Z[\sqrt5]$ is not: the golden ratio $\frac{1+\sqrt5}2$ satisfies $x^2-x-1=0$. The integrally closed rings of dimension one are the subject of Chapter~5.

\begin{lemma}\label{lem:intfield}
Let $A\subseteq B$ be domains, with $B$ integral over $A$. Then $B$ is a field if and only if $A$ is a field. Moreover every non-zero ideal of $B$ meets $A$ in a non-zero ideal.
\end{lemma}

\begin{proof}
Let $b\in B$, $b\ne0$, with a monic equation $b^n+a_{n-1}b^{n-1}+\dots+a_0=0$ of least degree. Then $a_0\ne0$: otherwise, as $B$ is a domain, we could divide by $b$. So $a_0=-b(b^{n-1}+\dots+a_1)$ is a non-zero element of $A\cap bB$, which proves
the last statement. If $A$ is a field, this formula shows $b^{-1}=-a_0^{-1}(b^{n-1}+\dots+a_1)\in B$. If $B$ is a field and $a\in A$ is non-zero, then $a^{-1}\in B$ is integral over $A$: $a^{-m}+c_{m-1}a^{-m+1}+\dots+c_0=0$, and multiplying by $a^{m-1}$
gives $a^{-1}=-(c_{m-1}+c_{m-2}a+\dots+c_0a^{m-1})\in A$.
\end{proof}

\begin{theorem}[Cohen and Seidenberg, 1946]\label{thm:goingup}
Let $B$ be integral over $A\subseteq B$.
\begin{enumerate}[label=(\alph*)]
\item (\emph{Incomparability}) If $\mathfrak q\subseteq\mathfrak q'$ are primes of $B$ with $\mathfrak q\cap A=\mathfrak q'\cap A$, then $\mathfrak q=\mathfrak q'$.
\item (\emph{Lying over}) For every prime $\mathfrak p$ of $A$ there is a prime $\mathfrak q$ of $B$ with $\mathfrak q\cap A=\mathfrak p$.
\item (\emph{Going up}) If $\mathfrak p\subseteq\mathfrak p'$ are primes of $A$ and $\mathfrak q\cap A=\mathfrak p$, there is a prime $\mathfrak q'\supseteq\mathfrak q$ with $\mathfrak q'\cap A=\mathfrak p'$.
\end{enumerate}
\end{theorem}

\begin{proof}
(a) Pass to $\bar A=A/(\mathfrak q\cap A)\subseteq\bar B=B/\mathfrak q$, domains with $\bar B$ integral over $\bar A$. The ideal $\mathfrak q'/\mathfrak q$ of $\bar B$ meets $\bar A$ in $0$, so it is $0$ by Lemma~\ref{lem:intfield}.

(b) Let $S=A\setminus\mathfrak p$. Then $B_{\mathfrak p}=S^{-1}B$ is integral over $A_{\mathfrak p}=S^{-1}A$ (divide a monic equation of $b$ by $s^n$), and $B_{\mathfrak p}\ne0$ since it contains $A_{\mathfrak p}\ne0$ (localisation is exact, so
$S^{-1}A\to S^{-1}B$ is injective). Let $\mathfrak n$ be a maximal ideal of $B_{\mathfrak p}$ (Zorn's lemma). By Lemma~\ref{lem:intfield} applied to $A_{\mathfrak p}/(\mathfrak n\cap A_{\mathfrak p})\subseteq B_{\mathfrak p}/\mathfrak n$, the prime
$\mathfrak n\cap A_{\mathfrak p}$ is maximal, so it is $\mathfrak pA_{\mathfrak p}$ by Corollary~\ref{cor:localring}. The inverse image $\mathfrak q$ of $\mathfrak n$ in $B$ is prime, and $\mathfrak q\cap A$ is the inverse image of
$\mathfrak pA_{\mathfrak p}$ in $A$, which is $\mathfrak p$ by Theorem~\ref{thm:locprimes}.

(c) Apply (b) to $A/\mathfrak p\subseteq B/\mathfrak q$, which is integral, and the prime $\mathfrak p'/\mathfrak p$; take for $\mathfrak q'$ the inverse image in $B$ of the prime found.
\end{proof}

\section{Krull dimension and Noether normalisation}

\begin{definition}\label{def:krulldim}
The \emph{Krull dimension} $\dim R$ is the supremum of the lengths $m$ of chains of prime ideals $\mathfrak p_0\subsetneq\mathfrak p_1\subsetneq\dots\subsetneq\mathfrak p_m$.
\end{definition}

Fields have dimension $0$, principal ideal domains that are not fields dimension $1$ (a chain is $0\subsetneq(p)$, maximal), and Theorem~\ref{thm:speczx} shows $\dim\Z[x]=2$, with chains like $0\subsetneq(x^2+1)\subsetneq(5,x-2)$.
Theorem~\ref{thm:goingup} gives at once:

\begin{corollary}\label{cor:dimintegral}
If $B$ is integral over $A\subseteq B$, then $\dim B=\dim A$.
\end{corollary}

\begin{proof}
A chain of primes in $B$ contracts to a chain in $A$, strictly increasing by incomparability: $\dim B\le\dim A$. A chain in $A$ lifts to a chain in $B$ by lying over and going up: $\dim A\le\dim B$.
\end{proof}

\begin{lemma}[Noether's change of variables]\label{lem:noethervar}
Let $K$ be a field and $f\in K[x_1,\dots,x_n]$ non-constant, $n\ge1$. For a large enough integer $r$, after the substitution $x_i=y_i+x_n^{r^i}$ for $i<n$, the polynomial $f$ becomes $c\,x_n^N+(\text{terms of lower degree in }x_n)$, with $c\in K$
non-zero and coefficients in $K[y_1,\dots,y_{n-1}]$.
\end{lemma}

\begin{proof}
Choose $r$ larger than every exponent occurring in $f$. A monomial $x^\alpha$ becomes $\prod_{i<n}(y_i+x_n^{r^i})^{\alpha_i}x_n^{\alpha_n}$, whose highest power of $x_n$ is $x_n^{e(\alpha)}$ with coefficient $1$, where
$e(\alpha)=\alpha_n+\alpha_1r+\alpha_2r^2+\dots+\alpha_{n-1}r^{n-1}$. Since all $\alpha_i<r$, the integers $e(\alpha)$ are distinct for distinct $\alpha$ (base-$r$ expansion). The monomial of $f$ with the largest $e(\alpha)$ gives the leading
term $c\,x_n^N$, with $c$ its coefficient and $N=e(\alpha)$, and no other monomial can cancel it.
\end{proof}

\begin{theorem}[Noether normalisation, 1926]\label{thm:noethernorm}
Let $A$ be a finitely generated algebra over a field $K$. There are elements $z_1,\dots,z_d\in A$, algebraically independent over $K$, such that $A$ is integral over $K[z_1,\dots,z_d]$.
\end{theorem}

\begin{proof}
Induction on the number $n$ of generators $a_1,\dots,a_n$ of $A$. If they are algebraically independent, take $z_i=a_i$, $d=n$ (this includes $n=0$). Otherwise there is a non-zero $f\in K[x_1,\dots,x_n]$ with $f(a_1,\dots,a_n)=0$; if $f$
were constant it would be $0$, so $f$ is non-constant. By Lemma~\ref{lem:noethervar}, with $b_i=a_i-a_n^{r^i}$ for $i<n$, the element $a_n$ satisfies an equation $c\,a_n^N+(\text{lower powers of }a_n)=0$ whose coefficients lie in $K[b_1,\dots,b_{n-1}]$; dividing by $c$, it is integral over
$A'=K[b_1,\dots,b_{n-1}]$. Then $A=A'[a_n]$ is integral over $A'$, since $a_i=b_i+a_n^{r^i}$. By induction $A'$ is integral over some $K[z_1,\dots,z_d]$ with the $z_i$ algebraically independent, and so is $A$ by
Corollary~\ref{cor:integralclosure}.
\end{proof}

Geometrically: every affine variety maps onto an affine space $K^d$ by a map with finite fibres. The number $d$ is its dimension, as the next theorem shows.

\begin{theorem}\label{thm:dimpoly}
For every field $K$, the polynomial ring $K[x_1,\dots,x_n]$ has dimension $n$. Consequently, if a finitely generated $K$-algebra $A$ is integral over $K[z_1,\dots,z_d]$ with the $z_i$ algebraically independent, then $\dim A=d$.
\end{theorem}

\begin{proof}
The chain $0\subsetneq(x_1)\subsetneq(x_1,x_2)\subsetneq\dots\subsetneq(x_1,\dots,x_n)$ consists of primes (each quotient is a polynomial ring), so $\dim\ge n$. We prove $\dim K[x_1,\dots,x_n]\le n$ by induction on $n$, the case $n=0$
being a field. Let $\mathfrak p_0\subsetneq\dots\subsetneq\mathfrak p_m$ be a chain with $m\ge1$; we may assume $\mathfrak p_0=0$, since replacing $\mathfrak p_0$ by $0$ keeps the chain strictly increasing. Choose a non-zero $f\in\mathfrak p_1$, which is not
constant since $\mathfrak p_1$ is proper. After the change of variables of Lemma~\ref{lem:noethervar}, which is an automorphism of $K[x_1,\dots,x_n]$ (its inverse is $y_i=x_i-x_n^{r^i}$), $f$ is $c$ times a monic polynomial in $x_n$
over $A'=K[y_1,\dots,y_{n-1}]$. Hence $B=K[x_1,\dots,x_n]/(f)$ is integral over the image of $A'$, and $A'\cap(f)=0$ because a non-zero multiple of $f$ has positive degree in $x_n$. So $B$ is integral over a copy of $A'$, and
$\dim B=\dim A'\le n-1$ by Corollary~\ref{cor:dimintegral} and induction. The chain $\mathfrak p_1/(f)\subsetneq\dots\subsetneq\mathfrak p_m/(f)$ in $B$ has length $m-1$, so $m-1\le n-1$.
For the last statement, combine with Corollary~\ref{cor:dimintegral}.
\end{proof}

\begin{beauty}
Mumford's drawing is not only a metaphor. In $\operatorname{Spec}\Z[x]$ the curve of $x^2+1$ is the spectrum of the Gaussian integers $\Z[i]$, and the way it meets the fibres is the way rational primes split in $\Z[i]$: in two
points for $p\equiv1\pmod4$, as Fermat's theorem on sums of two squares says, in one point otherwise, with the prime $2$ where the curve touches its fibre ``ramified''. A ring of integers of a number field is a curve over the
``line'' $\operatorname{Spec}\Z$, and much of modern number theory, from the Weil conjectures to the work of Arakelov, starts from taking that analogy literally. Chapter~5 develops the algebra of these arithmetic curves.
\end{beauty}

\section{Exercises}

\begin{exercise}[$\star$]\label{ex:subringsq}
Show that every subring of $\Q$ is a ring of fractions $S^{-1}\Z$, and that the ideals of $\Z_{(p)}$ are $0$ and the $p^k\Z_{(p)}$. What is $\operatorname{Spec}\Z_{(p)}$?
\end{exercise}

\begin{exercise}[$\star$]\label{ex:golden}
Show that $\Z[\sqrt5]$ is not integrally closed, and that $\Z[\frac{1+\sqrt5}2]$ is a finitely generated $\Z$-module. Find a monic equation over $\Z$ for $\sqrt2+\sqrt3$.
\end{exercise}

\begin{exercise}[$\star\star$]\label{ex:localiff}
Show that a ring is local if and only if its non-units form an ideal. Show that $R_f=0$ if and only if $f$ is nilpotent, and that the nilpotent elements of $R$ form the intersection of all its prime ideals.
\end{exercise}

\begin{exercise}[$\star\star$]\label{ex:specrx}
Describe $\operatorname{Spec}\R[x]$, and show that its closed points are in bijection with the upper half-plane $\{z\in\C:\operatorname{Im}z\ge0\}$.
\end{exercise}

\begin{exercise}[$\star\star$]\label{ex:quasicompact}
Show that $\operatorname{Spec}R$ is quasi-compact: every cover by open sets has a finite subcover. (The open sets $D(f)=\operatorname{Spec}R\setminus V(f)$ form a basis.)
\end{exercise}

\begin{exercise}[$\star\star$]\label{ex:hyperbola}
Let $A=\C[x,y]/(xy-1)$. Show that $A$ is not integral over $\C[x]$, but is integral over $\C[x+y]$. Interpret geometrically, with the projections of the hyperbola $xy=1$ to the two axes.
\end{exercise}

\section{Hints}

\begin{hint}{ex:subringsq}
Let $S$ be the set of denominators $b>0$ of reduced fractions $a/b$ in the subring; if $a/b$ is in it, so is $1/b$ (Bézout).
\end{hint}

\begin{hint}{ex:golden}
For the last part, square twice.
\end{hint}

\begin{hint}{ex:localiff}
For the last part, if $f$ is not nilpotent, $R_f\neq0$ has a prime ideal, and Theorem~\ref{thm:locprimes} pulls it back.
\end{hint}

\begin{hint}{ex:specrx}
$\R[x]$ is a principal ideal domain; its monic irreducibles have degree $1$ or $2$.
\end{hint}

\begin{hint}{ex:quasicompact}
A cover by $D(f_\lambda)$ means that no prime contains all the $f_\lambda$, so they generate the unit ideal; and $1$ is a finite combination.
\end{hint}

\begin{hint}{ex:hyperbola}
Is $y$ integral over $\C[x]$? Compute $y^2-(x+y)y+1$.
\end{hint}

\section{Solutions}

\begin{solution}{ex:subringsq}
\idea A subring of $\Q$ is determined by which denominators it allows.
\writeup Let $R\subseteq\Q$ be a subring and $S$ the set of positive integers $b$ such that $1/b\in R$, a multiplicative set. If $a/b\in R$ in lowest terms, write $ua+vb=1$; then $\frac1b=u\frac ab+v\in R$, so $b\in S$. Hence $R=S^{-1}\Z$.
The ideals of $\Z_{(p)}$: a non-zero element is $p^k\cdot u$ with $u$ a unit (numerator and denominator prime to $p$), so a non-zero ideal is generated by its element with the least $k$: it is $p^k\Z_{(p)}$. The primes are $0$ and
$p\Z_{(p)}$: $\operatorname{Spec}\Z_{(p)}$ has two points, a closed point and a generic point, the prime $p$ seen up close.
\end{solution}

\begin{solution}{ex:golden}
\idea An integral element of the fraction field outside the ring proves it is not integrally closed.
\writeup $\varphi=\frac{1+\sqrt5}2$ is a root of $x^2-x-1$, lies in $\Q(\sqrt5)$, the field of fractions of $\Z[\sqrt5]$, and is not in $\Z[\sqrt5]$ (its coefficient of $\sqrt5$ is $\frac12$). The ring $\Z[\varphi]=\Z+\Z\varphi$ is
generated by $1,\varphi$ since $\varphi^2=\varphi+1$ (Proposition~\ref{prop:integralcrit}). For $t=\sqrt2+\sqrt3$: $t^2=5+2\sqrt6$, $(t^2-5)^2=24$, so $t^4-10t^2+1=0$.
\end{solution}

\begin{solution}{ex:localiff}
\idea In a local ring the non-units are $\mathfrak m$; conversely an ideal of non-units contains every proper ideal.
\writeup If $R$ is local, the non-units are the elements of $\mathfrak m$ (they generate proper ideals, contained in a maximal ideal), an ideal. Conversely, if the non-units form an ideal $\mathfrak m$, every proper ideal consists of non-units, so lies in
$\mathfrak m$, which is then the only maximal ideal. Next, $R_f=0$ iff $1/1=0$ in $R_f$ iff $f^n\cdot1=0$ for some $n$. Finally, a nilpotent $f$ lies in every prime, since $f^n=0\in\mathfrak p$. If $f$ is not nilpotent, $R_f\ne0$ has a maximal
ideal (Zorn), whose inverse image in $R$ is a prime not containing $f$ by Theorem~\ref{thm:locprimes}.
\end{solution}

\begin{solution}{ex:specrx}
\idea The primes of a principal ideal domain are $0$ and the irreducibles.
\writeup $\operatorname{Spec}\R[x]$ consists of the generic point $(0)$, the closed points $(x-a)$ for $a\in\R$, and the closed points $(x^2+bx+c)$ with $b^2<4c$. The latter correspond to pairs of conjugate non-real roots $z,\bar z$, and so
to the $z$ with $\operatorname{Im}z>0$; the former to the real $a$, the boundary of the half-plane. So the closed points form $\{\operatorname{Im}z\ge0\}$, which is $\C$ folded along the real axis: $\operatorname{Spec}\R[x]$ is
$\operatorname{Spec}\C[x]$ divided by complex conjugation.
\end{solution}

\begin{solution}{ex:quasicompact}
\idea Covering $\operatorname{Spec}R$ by open sets is an algebraic statement about generating the unit ideal.
\writeup Every open set is a union of basic sets $D(f)$, because $\operatorname{Spec}R\setminus V(I)=\bigcup_{f\in I}D(f)$; so it is enough to treat a cover by sets $D(f_\lambda)$. Then no prime contains all the $f_\lambda$, so the ideal they
generate is contained in no maximal ideal, hence is $R$ (Zorn's lemma, Volume~6, Theorem~1.19). Then $1=\sum_{i=1}^kr_if_{\lambda_i}$, a finite sum, and the $D(f_{\lambda_i})$ already cover: a prime containing all $f_{\lambda_i}$ would contain $1$.
\end{solution}

\begin{solution}{ex:hyperbola}
\idea Over $\C[x]$, the function $y=1/x$ has a pole; over $\C[x+y]$, $y$ is a root of a monic quadratic.
\writeup If $y$ were integral over $\C[x]$, then $y^n+p_{n-1}(x)y^{n-1}+\dots+p_0(x)=0$ in $A$. In $A$, $x$ is invertible with inverse $y$, and $A\cong\C[x,x^{-1}]$; multiplying by $x^n$ gives $1+p_{n-1}(x)x+\dots+p_0(x)x^n=0$ in
$\C[x,x^{-1}]$, impossible since the left side has constant term $1$ and is a polynomial in $x$. But $y^2-(x+y)y+1=1-xy=0$ in $A$: $y$ is integral over $\C[x+y]$, and so is $x=(x+y)-y$, hence $A$. Geometrically, the projection
$(x,y)\mapsto x$ from the hyperbola misses the point $0$, whose fibre is ``at infinity'', while $(x,y)\mapsto x+y$ is onto $\C$ with two points in almost every fibre, and with finite fibres everywhere: it is a Noether normalisation.
\end{solution}

\section*{Sources and further reading}
\addcontentsline{toc}{section}{Sources and further reading}

E.~Noether, ``Idealtheorie in Ringbereichen'', \emph{Math.\ Ann.} 83 (1921), 24--66, and ``Der Endlichkeitssatz der Invarianten endlicher linearer Gruppen der Charakteristik $p$'', \emph{Nachr.\ Ges.\ Wiss.\ Göttingen} (1926), 28--35, where
the normalisation lemma appears; W.~Krull, ``Dimensionstheorie in Stellenringen'', \emph{J.\ Reine Angew.\ Math.} 179 (1938), 204--226; I.~S.~Cohen and A.~Seidenberg, ``Prime ideals and integral dependence'', \emph{Bull.\ Amer.\ Math.\ Soc.} 52
(1946), 252--261. The picture of $\operatorname{Spec}\Z[x]$ is in D.~Mumford, \emph{The Red Book of Varieties and Schemes} (Springer Lecture Notes in Mathematics 1358, 1988), first circulated as mimeographed notes in 1967.

For more: M.~F.~Atiyah and I.~G.~Macdonald, \emph{Introduction to Commutative Algebra} (1969), Chapters~3, 5 and 11; D.~Eisenbud, \emph{Commutative Algebra with a View Toward Algebraic Geometry} (1995), Chapters~2, 4, 8 and 13;
R.~Vakil, \emph{The Rising Sea: Foundations of Algebraic Geometry}, lecture notes revised many times since 2010, Chapters~3 and 4, for the spectrum and many more pictures.
